\documentclass[12pt]{article}

%------------- MH5200-STYLE PREAMBLE (compatible subset) -------------
\usepackage[margin=1in]{geometry}
\usepackage{amsmath,amssymb,mathtools}
\usepackage{enumitem}
\usepackage{bm}
\usepackage{hyperref}
\usepackage{fancyhdr}
\usepackage{draftwatermark}
\usepackage{xcolor}

%------------- TITLE AND METADATA (REQUIRED STRUCTURE) -------------
\newcommand{\CourseCode}{MH1201}
\newcommand{\CourseName}{Linear Algebra II}
\newcommand{\DocType}{Tutorial 10 (Week 12 \& 13) -- Problems \& Solutions}


\title{
\CourseCode\ \CourseName \\[0.3em]
\large \DocType
}
\author{
  Academic Year 2025/2026, Semester 2 \\[0.25em]
  \textit{Quantitative Research Society @NTU}
}
\date{\today}

%------------- HEADER & FOOTER (for page 2 onward) -------------
\fancypagestyle{main}{
  \fancyhf{}
  \fancyhead[L]{\CourseCode\ \CourseName}
  \fancyhead[R]{\href{https://qrsntu.org}{QRS@NTU}}
  \fancyfoot[C]{\thepage}
  \renewcommand{\headrulewidth}{0.4pt}
  \renewcommand{\footrulewidth}{0pt}
}

%------------- SIMPLE FOOTER (page 1 only) -------------
\fancypagestyle{plainfirst}{
  \fancyhf{}
  \renewcommand{\headrulewidth}{0pt}
  \renewcommand{\footrulewidth}{0pt}
  \fancyfoot[C]{\scriptsize\textcolor{gray}{\textit{For academic reference only. This document is provided for educational purposes and does not constitute an official question paper, answer key, or any formally authorised academic material. Its content is not guaranteed to be complete or accurate.}}}
}

%------------- WATERMARK (MH5200-style approach) -------------
\SetWatermarkText{Unofficial — QRS@NTU — qrsntu.org}
\SetWatermarkScale{1}
\SetWatermarkLightness{0.99}

%------------- COMMON MACROS -------------
\newcommand{\df}{\coloneqq}
\newcommand{\R}{\mathbb{R}}
\newcommand{\C}{\mathbb{C}}
\newcommand{\cL}{\mathcal{L}}
\newcommand{\cP}{\mathcal{P}}
\DeclareMathOperator{\Span}{span}
\DeclareMathOperator{\Rank}{rank}
\DeclareMathOperator{\Nullity}{nullity}
\DeclareMathOperator{\Range}{range}
\DeclareMathOperator{\tr}{trace}
\DeclareMathOperator{\proj}{proj}
\DeclareMathOperator{\diag}{diag}
\newcommand{\Id}{\mathrm{id}}
\newcommand{\ip}[2]{\left\langle #1,#2\right\rangle}
%----------------------------------------
% Optional: tighter list defaults (feel free to adjust)
\setlist[enumerate]{itemsep=0.3em, topsep=0.3em}
\setlist[itemize]{itemsep=0.2em, topsep=0.2em}

\setcounter{secnumdepth}{-1} % Globally removes numbers from all sections
\setcounter{tocdepth}{3}     % Ensures \subsubsection level shows in TOC

\begin{document}

%------------- COVER PAGE -------------
\maketitle
\thispagestyle{plainfirst}
\pagestyle{main}

%====================================================
% OVERVIEW (PAGE 1)
%====================================================
\begin{center}
  \rule{0.8\textwidth}{0.4pt}\\[4pt]
  \textbf{Overview}\\[2pt]
  \rule{0.8\textwidth}{0.4pt}
\end{center}

\noindent
This tutorial emphasizes:
\begin{itemize}[leftmargin=*]
  \item Gram--Schmidt in abstract inner product spaces and stability on orthogonal inputs;
  \item producing orthonormal bases that preserve a $T$-invariant flag, hence preserving upper-triangularity;
  \item orthogonal polynomials on $[-1,1]$ and Fourier coefficients in an orthonormal basis;
  \item inner products defined by Hermitian positive definite matrices and complex Gram--Schmidt;
  \item orthogonal complements: $W\subseteq W^{\perp\perp}$ always, and $W=W^{\perp\perp}$ in finite dimensions;
  \item explicit orthogonal decomposition $x=a+b$ with $a\in V$, $b\in V^\perp$ in $\R^4$;
  \item Frobenius inner product on $\R^{n,n}$ and describing $D^\perp$ for diagonal $D$.
\end{itemize}

\newpage
\tableofcontents
\newpage

%====================================================
% QUESTION 1
%====================================================
\section{Problem 1 (Gram--Schmidt on an orthogonal basis)}
\subsection{Problem}
Let $(V,\ip{\cdot}{\cdot})$ be a finite-dimensional inner product space, with $\dim(V)=n\ge 1$.
Let $\beta=(b_1,\dots,b_n)$ be an orthogonal ordered basis for $(V,\ip{\cdot}{\cdot})$.
Show that applying Gram--Schmidt to $\beta$ with respect to $(V,\ip{\cdot}{\cdot})$ yields $\beta$ back.

\subsection{Solution}

\subsubsection{Method 1: Projections vanish because of orthogonality}
Recall the (non-normalized) Gram--Schmidt construction:
\[
u_1=b_1,\qquad
u_k \df b_k-\sum_{j=1}^{k-1}\proj_{u_j}(b_k),
\quad
\proj_{u_j}(b_k)\df \frac{\ip{b_k}{u_j}}{\ip{u_j}{u_j}}u_j,\qquad 2\le k\le n.
\]
We prove by induction that $u_k=b_k$ for all $k$.

For $k=1$, $u_1=b_1$ is immediate.

Assume $u_j=b_j$ for all $j<k$. Since $\beta$ is orthogonal, for each $j<k$,
\[
\ip{b_k}{u_j}=\ip{b_k}{b_j}=0.
\]
Hence each projection term is $0$, so
\[
u_k=b_k-\sum_{j=1}^{k-1} 0=b_k.
\]
Thus Gram--Schmidt returns the same ordered orthogonal basis $\beta$.

\[
\boxed{\,\text{If }\beta\text{ is already orthogonal, Gram--Schmidt outputs }\beta\text{ unchanged.}\,}
\]

\subsubsection{Method 2: Uniqueness from nested spans}
In Gram--Schmidt, the output vectors satisfy:
\[
u_k\in \Span(b_1,\dots,b_k),
\qquad
\Span(u_1,\dots,u_k)=\Span(b_1,\dots,b_k),
\qquad
\ip{u_k}{u_j}=0\ (j<k).
\]
Because $\beta$ is orthogonal, the \emph{unique} vector in $\Span(b_1,\dots,b_k)$ which is orthogonal to
$\Span(b_1,\dots,b_{k-1})$ and has nonzero component along $b_k$ is a scalar multiple of $b_k$.
But Gram--Schmidt uses the formula $u_k=b_k-(\text{component in }\Span(u_1,\dots,u_{k-1}))$, which preserves the $b_k$-component,
so the scalar multiple must be $1$, hence $u_k=b_k$.

%====================================================
% QUESTION 2
%====================================================
\newpage
\section{Problem 2 (Orthonormal basis preserving upper-triangularity)}
\subsection{Problem}
Let
\[
\beta \df ((1,0,0),(1,1,1),(1,1,2)),
\]
which is an ordered basis for $\R^3$ (with the standard inner product $\cdot$).
Let $T\in \cL(\R^3)$, and suppose that $[T]_{\beta}^{\beta}\in \R^{3,3}$ is upper triangular.
Derive an orthonormal ordered basis $\gamma$ for $(\R^3,\cdot)$ for which $[T]_{\gamma}^{\gamma}\in\R^{3,3}$ is upper triangular.

\subsection{Solution}

\subsubsection{Method 1: Gram--Schmidt preserves the flag; hence upper triangular}
Write $\beta=(b_1,b_2,b_3)$ with
\[
b_1=(1,0,0),\quad b_2=(1,1,1),\quad b_3=(1,1,2).
\]
Since $[T]_{\beta}^{\beta}$ is upper triangular, we have the equivalent $T$-invariance of the coordinate flag:
\[
T\bigl(\Span(b_1)\bigr)\subseteq \Span(b_1),\quad
T\bigl(\Span(b_1,b_2)\bigr)\subseteq \Span(b_1,b_2),\quad
T\bigl(\Span(b_1,b_2,b_3)\bigr)\subseteq \Span(b_1,b_2,b_3)=\R^3.
\]
Apply Gram--Schmidt to $\beta$ in the standard inner product.

\medskip
\noindent\textbf{Step 1: Compute $\gamma$ by Gram--Schmidt.}
Let
\[
u_1=b_1,\qquad
u_2=b_2-\proj_{u_1}(b_2),\qquad
u_3=b_3-\proj_{u_1}(b_3)-\proj_{u_2}(b_3).
\]
Compute:
\[
\proj_{u_1}(b_2)=\frac{b_2\cdot u_1}{u_1\cdot u_1}u_1=\frac{1}{1}b_1=b_1,
\quad\Rightarrow\quad u_2=b_2-b_1=(0,1,1).
\]
Similarly,
\[
\proj_{u_1}(b_3)=\frac{b_3\cdot u_1}{u_1\cdot u_1}u_1=b_1,
\quad\Rightarrow\quad b_3-\proj_{u_1}(b_3)=(0,1,2).
\]
Now
\[
\proj_{u_2}(b_3)=\proj_{u_2}\bigl((0,1,2)\bigr)
=\frac{(0,1,2)\cdot (0,1,1)}{(0,1,1)\cdot(0,1,1)}(0,1,1)
=\frac{3}{2}(0,1,1).
\]
Thus
\[
u_3=(0,1,2)-\frac{3}{2}(0,1,1)=\left(0,-\frac12,\frac12\right).
\]
Normalize:
\[
\|u_1\|=1,\quad \|u_2\|=\sqrt2,\quad \|u_3\|=\sqrt{\frac12}.
\]
Hence an orthonormal basis is
\[
\gamma=(g_1,g_2,g_3)\df
\left(
(1,0,0),\ \left(0,\frac{1}{\sqrt2},\frac{1}{\sqrt2}\right),\ \left(0,-\frac{1}{\sqrt2},\frac{1}{\sqrt2}\right)
\right).
\]
\[
\boxed{\,
\gamma=\left(
(1,0,0),\ \left(0,\frac{1}{\sqrt2},\frac{1}{\sqrt2}\right),\ \left(0,-\frac{1}{\sqrt2},\frac{1}{\sqrt2}\right)
\right)\ \text{is orthonormal.}
\,}
\]

\medskip
\noindent\textbf{Step 2: Show $[T]_\gamma^\gamma$ is upper triangular.}
A key Gram--Schmidt property is
\[
g_k\in \Span(b_1,\dots,b_k)\qquad (k=1,2,3).
\]
Therefore
\[
\Span(g_1,\dots,g_k)=\Span(b_1,\dots,b_k)\qquad (k=1,2,3).
\]
Since $T(\Span(b_1,\dots,b_k))\subseteq \Span(b_1,\dots,b_k)$, we get
\[
T(\Span(g_1,\dots,g_k))\subseteq \Span(g_1,\dots,g_k)\qquad (k=1,2,3).
\]
This is exactly the condition that the matrix of $T$ in the ordered basis $\gamma$ is upper triangular.

\[
\boxed{\, [T]_{\gamma}^{\gamma}\ \text{is upper triangular.}\,}
\]

\subsubsection{Method 2: Upper-triangular similarity via an upper-triangular change-of-basis matrix}
Let $B=[T]_{\beta}^{\beta}$ (upper triangular). Let $P$ be the change-of-basis matrix from $\gamma$-coordinates to $\beta$-coordinates:
\[
P \df [\Id]_{\beta}^{\gamma}.
\]
Because Gram--Schmidt produces $g_k\in\Span(b_1,\dots,b_k)$, the coordinate vector $[g_k]_\beta$ has no components in $b_{k+1},\dots,b_3$,
so $P$ is upper triangular with nonzero diagonal (hence invertible).
Then
\[
[T]_{\gamma}^{\gamma}=P^{-1}\,[T]_{\beta}^{\beta}\,P=P^{-1}BP.
\]
A product (and inverse) of upper-triangular invertible matrices is upper triangular, hence $[T]_{\gamma}^{\gamma}$ is upper triangular.

%====================================================
% QUESTION 3
%====================================================
\newpage
\section{Problem 3 (Orthogonal polynomials on $[-1,1]$)}
\subsection{Problem}
Define a function $\ip{\cdot}{\cdot}:\cP_3(\R)\times \cP_3(\R)\to \R$ as follows:
\[
\forall P,Q\in\cP_3(\R):\qquad
\ip{P}{Q}\df \int_{-1}^{1} P(x)\,Q(x)\,dx.
\]
\begin{enumerate}[label=(\alph*)]
\item Show that $\ip{\cdot}{\cdot}$ is an inner product on $\cP_3(\R)$.
\item Apply Gram--Schmidt to $(1,X,X^2,X^3)$ with respect to $(\cP_3(\R),\ip{\cdot}{\cdot})$.
\item Normalize your answer to Part (b) to obtain an orthonormal basis $\beta$ for $(\cP_3(\R),\ip{\cdot}{\cdot})$.
\item Compute the $\beta$-Fourier coefficients of $1-X+X^2-X^3$ with respect to $(\cP_3(\R),\ip{\cdot}{\cdot})$.
\end{enumerate}

\subsection{Solution}

\subsubsection{Method 1: Direct axiom verification + explicit Gram--Schmidt computations}
\begin{enumerate}[label=(\alph*)]
\item \textbf{Inner product axioms.}
Bilinearity follows from linearity of the integral:
\[
\ip{aP_1+bP_2}{Q}=a\ip{P_1}{Q}+b\ip{P_2}{Q},\qquad
\ip{P}{aQ_1+bQ_2}=a\ip{P}{Q_1}+b\ip{P}{Q_2}.
\]
Symmetry holds because $P(x)Q(x)=Q(x)P(x)$:
\[
\ip{P}{Q}=\int_{-1}^1 P Q=\int_{-1}^1 Q P=\ip{Q}{P}.
\]
Positive definiteness: for $P\neq 0$,
\[
\ip{P}{P}=\int_{-1}^1 P(x)^2\,dx>0.
\]
Indeed, $P(x)^2\ge 0$ and is continuous; if the integral were $0$, then $P(x)^2=0$ for all $x\in[-1,1]$, so $P$ vanishes on an interval,
hence $P$ is the zero polynomial, contradiction. Therefore $\ip{\cdot}{\cdot}$ is an inner product.

\[
\boxed{\,\ip{\cdot}{\cdot}\text{ is an inner product on }\cP_3(\R).\,}
\]

\item \textbf{Gram--Schmidt on $(1,X,X^2,X^3)$.}
Let $v_0=1$, $v_1=X$, $v_2=X^2$, $v_3=X^3$.
Define
\[
u_0=v_0,\qquad
u_1=v_1-\proj_{u_0}(v_1),\qquad
u_2=v_2-\proj_{u_0}(v_2)-\proj_{u_1}(v_2),\qquad
u_3=v_3-\sum_{j=0}^{2}\proj_{u_j}(v_3).
\]

Compute required inner products using parity:
\[
\int_{-1}^1 x\,dx=0,\quad \int_{-1}^1 x^3\,dx=0,\quad \int_{-1}^1 x^5\,dx=0.
\]
Also
\[
\int_{-1}^1 1\,dx=2,\quad \int_{-1}^1 x^2\,dx=\frac{2}{3},\quad \int_{-1}^1 x^4\,dx=\frac{2}{5},\quad \int_{-1}^1 x^6\,dx=\frac{2}{7}.
\]

\underline{Step 1: $u_0=1$.}

\underline{Step 2: $u_1$.}
Since $\ip{X}{1}=\int_{-1}^1 x\,dx=0$, we get
\[
u_1=X.
\]

\underline{Step 3: $u_2$.}
Compute
\[
\proj_{u_0}(X^2)=\frac{\ip{X^2}{1}}{\ip{1}{1}}1=\frac{\frac{2}{3}}{2}\cdot 1=\frac{1}{3}.
\]
Also $\ip{X^2}{X}=\int_{-1}^1 x^3\,dx=0$, so $\proj_{u_1}(X^2)=0$.
Hence
\[
u_2=X^2-\frac{1}{3}.
\]

\underline{Step 4: $u_3$.}
We have $\ip{X^3}{1}=\int_{-1}^1 x^3\,dx=0$, and
\[
\proj_{u_1}(X^3)=\frac{\ip{X^3}{X}}{\ip{X}{X}}X
=\frac{\int_{-1}^1 x^4\,dx}{\int_{-1}^1 x^2\,dx}X
=\frac{\frac{2}{5}}{\frac{2}{3}}X=\frac{3}{5}X.
\]
Finally,
\[
\ip{X^3}{u_2}=\int_{-1}^1 x^3\left(x^2-\frac13\right)\,dx
=\int_{-1}^1 \left(x^5-\frac13 x^3\right)\,dx=0,
\]
so $\proj_{u_2}(X^3)=0$. Therefore
\[
u_3=X^3-\frac{3}{5}X.
\]

Thus Gram--Schmidt yields the orthogonal list
\[
\boxed{\,\bigl(1,\ X,\ X^2-\tfrac13,\ X^3-\tfrac35 X\bigr).\,}
\]

\item \textbf{Normalize to an orthonormal basis $\beta$.}
Compute norms:

\underline{$u_0=1$:} $\|u_0\|^2=\int_{-1}^1 1\,dx=2$, so $\|u_0\|=\sqrt2$.

\underline{$u_1=X$:} $\|u_1\|^2=\int_{-1}^1 x^2\,dx=\frac{2}{3}$, so $\|u_1\|=\sqrt{\frac{2}{3}}$.

\underline{$u_2=X^2-\frac13$:}
\[
\|u_2\|^2=\int_{-1}^1 \left(x^2-\frac13\right)^2dx
=\int_{-1}^1\left(x^4-\frac{2}{3}x^2+\frac{1}{9}\right)dx
=\frac{2}{5}-\frac{2}{3}\cdot\frac{2}{3}+\frac{1}{9}\cdot 2
=\frac{8}{45}.
\]
So $\|u_2\|=\sqrt{\frac{8}{45}}$.

\underline{$u_3=X^3-\frac35 X$:}
\[
\|u_3\|^2=\int_{-1}^1\left(x^3-\frac{3}{5}x\right)^2dx
=\int_{-1}^1\left(x^6-\frac{6}{5}x^4+\frac{9}{25}x^2\right)dx
=\frac{2}{7}-\frac{6}{5}\cdot\frac{2}{5}+\frac{9}{25}\cdot\frac{2}{3}
=\frac{8}{175}.
\]
So $\|u_3\|=\sqrt{\frac{8}{175}}$.

Define the orthonormal basis $\beta=(\beta_0,\beta_1,\beta_2,\beta_3)$ by $\beta_j=u_j/\|u_j\|$:
\[
\beta_0=\frac{1}{\sqrt2},\qquad
\beta_1=\sqrt{\frac{3}{2}}\,X,\qquad
\beta_2=\sqrt{\frac{45}{8}}\left(X^2-\frac13\right),\qquad
\beta_3=\sqrt{\frac{175}{8}}\left(X^3-\frac35 X\right).
\]
\[
\boxed{\,\beta=\left(\frac{1}{\sqrt2},\ \sqrt{\frac{3}{2}}X,\ \sqrt{\frac{45}{8}}\left(X^2-\frac13\right),\ \sqrt{\frac{175}{8}}\left(X^3-\frac35 X\right)\right)\ \text{is orthonormal.}\,}
\]

\item \textbf{$\beta$-Fourier coefficients of $f(X)=1-X+X^2-X^3$.}
Since $\beta$ is orthonormal, the Fourier coefficient of $f$ along $\beta_j$ is
\[
c_j=\ip{f}{\beta_j}.
\]

\underline{$c_0$:}
\[
c_0=\ip{f}{\tfrac{1}{\sqrt2}}=\frac{1}{\sqrt2}\int_{-1}^1 (1-x+x^2-x^3)\,dx
=\frac{1}{\sqrt2}\left(2+\frac{2}{3}\right)=\frac{8}{3\sqrt2}.
\]

\underline{$c_1$:}
\[
c_1=\ip{f}{\sqrt{\frac32}x}=\sqrt{\frac32}\int_{-1}^1 (1-x+x^2-x^3)x\,dx
=\sqrt{\frac32}\int_{-1}^1(x-x^2+x^3-x^4)\,dx
=\sqrt{\frac32}\left(-\frac{2}{3}-\frac{2}{5}\right)
=-\frac{16}{15}\sqrt{\frac32}.
\]

\underline{$c_2$:}
\[
c_2=\ip{f}{\sqrt{\frac{45}{8}}\left(x^2-\frac13\right)}
=\sqrt{\frac{45}{8}}\int_{-1}^1 (1-x+x^2-x^3)\left(x^2-\frac13\right)\,dx.
\]
Expanding and using parity (odd terms integrate to $0$) gives
\[
\int_{-1}^1 (1-x+x^2-x^3)\left(x^2-\frac13\right)\,dx
=\frac{8}{45}.
\]
Thus
\[
c_2=\sqrt{\frac{45}{8}}\cdot \frac{8}{45}=\frac{4\sqrt5}{15\sqrt2}.
\]

\underline{$c_3$:}
\[
c_3=\ip{f}{\sqrt{\frac{175}{8}}\left(x^3-\frac35 x\right)}
=\sqrt{\frac{175}{8}}\int_{-1}^1 (1-x+x^2-x^3)\left(x^3-\frac35 x\right)\,dx.
\]
Keeping only even powers after expansion yields
\[
\int_{-1}^1 (1-x+x^2-x^3)\left(x^3-\frac35 x\right)\,dx
=-\frac{8}{175}.
\]
Hence
\[
c_3=\sqrt{\frac{175}{8}}\cdot\left(-\frac{8}{175}\right)=-\frac{4\sqrt7}{35\sqrt2}.
\]

Therefore the $\beta$-Fourier coefficients are
\[
\boxed{\,
(c_0,c_1,c_2,c_3)=\left(\frac{8}{3\sqrt2},\ -\frac{16}{15}\sqrt{\frac32},\ \frac{4\sqrt5}{15\sqrt2},\ -\frac{4\sqrt7}{35\sqrt2}\right).
\,}
\]
\end{enumerate}

\subsubsection{Method 2: Recognize (scaled) Legendre polynomials}
The Gram--Schmidt output in (b) is (up to scaling) the Legendre family on $[-1,1]$:
\[
P_0=1,\quad P_1=x,\quad P_2=x^2-\frac13,\quad P_3=x^3-\frac35 x,
\]
which are mutually orthogonal under $\int_{-1}^1(\cdot)(\cdot)$.
Normalization constants follow from explicit integrals $\int_{-1}^1 P_k(x)^2dx$.
Fourier coefficients are $\ip{f}{\beta_k}$ (orthonormality), matching Method 1.

%====================================================
% QUESTION 4
%====================================================
\newpage
\section{Problem 4 (Complex inner product via a Hermitian positive definite matrix)}
\subsection{Problem}
Define a function $\ip{\cdot}{\cdot}:\C^{3,1}\times \C^{3,1}\to \C$ as follows:
\[
\forall u,v\in \C^{3,1}:\qquad
\ip{u}{v}\df u^{\dagger}
\begin{bmatrix}
2 & 1+i & 0\\
1-i & 3 & 2\\
0 & 2 & 4
\end{bmatrix}
v,
\]
where ${}^\dagger$ denotes conjugate transposition.
\begin{enumerate}[label=(\alph*)]
\item Show that $\ip{\cdot}{\cdot}$ is an inner product on $\C^{3,1}$.
\item Apply Gram--Schmidt to the standard ordered basis for $\C^{3,1}$ with respect to $(\C^{3,1},\ip{\cdot}{\cdot})$.
\item Normalize your answer to Part (b) to obtain an orthonormal basis $\beta$ for $(\C^{3,1},\ip{\cdot}{\cdot})$.
\item Compute the $\beta$-Fourier coefficients of $\begin{bmatrix}1\\0\\ i\end{bmatrix}$ with respect to $(\C^{3,1},\ip{\cdot}{\cdot})$.
\end{enumerate}

\subsection{Solution}
Let
\[
H\df
\begin{bmatrix}
2 & 1+i & 0\\
1-i & 3 & 2\\
0 & 2 & 4
\end{bmatrix},
\qquad
\ip{u}{v}=u^\dagger H v.
\]
We use the standard convention: conjugate-linear in the first argument and linear in the second.

\subsubsection{Method 1: Hermitian positive definite $\Rightarrow$ inner product; then compute Gram--Schmidt explicitly}
\begin{enumerate}[label=(\alph*)]
\item \textbf{Axiom 1: sesquilinearity.}
For $a,b\in\C$ and $u_1,u_2,v\in\C^{3,1}$,
\[
\ip{au_1+bu_2}{v}=(au_1+bu_2)^\dagger H v=\overline{a}\,u_1^\dagger Hv+\overline{b}\,u_2^\dagger Hv
=\overline{a}\ip{u_1}{v}+\overline{b}\ip{u_2}{v},
\]
and similarly $\ip{u}{av_1+bv_2}=a\ip{u}{v_1}+b\ip{u}{v_2}$.

\smallskip
\noindent\textbf{Axiom 2: conjugate symmetry.}
Since $H^\dagger=H$ (check immediately from entries), we have
\[
\overline{\ip{v}{u}}=\overline{v^\dagger H u}=u^\dagger H^\dagger v=u^\dagger H v=\ip{u}{v}.
\]

\smallskip
\noindent\textbf{Axiom 3: positive definiteness.}
Because $H$ is Hermitian, $\ip{u}{u}=u^\dagger H u\in\R$.
We show $u^\dagger H u>0$ for $u\neq 0$ using Sylvester's criterion (leading principal minors):
\[
\det([2])=2>0,\qquad
\det\begin{bmatrix}2&1+i\\1-i&3\end{bmatrix}
=6-(1+i)(1-i)=6-2=4>0,
\]
and
\[
\det(H)=2\det\begin{bmatrix}3&2\\2&4\end{bmatrix}-(1+i)\det\begin{bmatrix}1-i&2\\0&4\end{bmatrix}
=2(12-4)- (1+i)\cdot 4(1-i)
=16-4\cdot 2=8>0.
\]
Hence $H$ is Hermitian positive definite, so $u^\dagger H u>0$ for all $u\neq 0$.
Therefore $\ip{\cdot}{\cdot}$ is an inner product.

\[
\boxed{\,\ip{\cdot}{\cdot}\text{ is an inner product on }\C^{3,1}. \,}
\]

\item \textbf{Gram--Schmidt on the standard basis.}
Let $e_1,e_2,e_3$ be the standard basis of $\C^{3,1}$.
Define
\[
u_1=e_1,\qquad
u_2=e_2-\proj_{u_1}(e_2),\qquad
u_3=e_3-\proj_{u_1}(e_3)-\proj_{u_2}(e_3),
\]
where $\proj_{u}(v)=\dfrac{\ip{u}{v}}{\ip{u}{u}}u$.

Compute the needed inner products using $\ip{e_i}{e_j}=H_{ij}$:
\[
\ip{u_1}{u_1}=\ip{e_1}{e_1}=2,\qquad
\ip{u_1}{e_2}=\ip{e_1}{e_2}=1+i,\qquad
\ip{u_1}{e_3}=\ip{e_1}{e_3}=0.
\]
Hence
\[
u_2=e_2-\frac{1+i}{2}e_1.
\]
Next,
\[
\ip{u_2}{u_2}=2,
\qquad
\ip{u_2}{e_3}=2,
\]
so $\proj_{u_2}(e_3)=\dfrac{2}{2}u_2=u_2$ and $\proj_{u_1}(e_3)=0$.
Thus
\[
u_3=e_3-u_2=e_3-e_2+\frac{1+i}{2}e_1.
\]
So Gram--Schmidt yields the orthogonal ordered basis
\[
\boxed{\,
(u_1,u_2,u_3)=\left(e_1,\ e_2-\frac{1+i}{2}e_1,\ \frac{1+i}{2}e_1-e_2+e_3\right).
\,}
\]

\item \textbf{Normalize to an orthonormal basis $\beta$.}
We already computed $\ip{u_1}{u_1}=2$ and $\ip{u_2}{u_2}=2$.
Also,
\[
\ip{u_3}{u_3}=2
\]
(one direct check is $Hu_3=(0,0,2)^{\mathsf T}$, hence $u_3^\dagger Hu_3=2$).
Therefore $\|u_k\|=\sqrt2$ for $k=1,2,3$.

Define
\[
\beta=(\beta_1,\beta_2,\beta_3)\df \left(\frac{u_1}{\sqrt2},\ \frac{u_2}{\sqrt2},\ \frac{u_3}{\sqrt2}\right).
\]
Explicitly,
\[
\boxed{
\beta=
\left(
\frac{1}{\sqrt2}\!\begin{bmatrix}1\\0\\0\end{bmatrix},\ 
\frac{1}{\sqrt2}\!\begin{bmatrix}-\frac{1+i}{2}\\1\\0\end{bmatrix},\ 
\frac{1}{\sqrt2}\!\begin{bmatrix}\frac{1+i}{2}\\-1\\1\end{bmatrix}
\right)
\ \text{is orthonormal in }(\C^{3,1},\ip{\cdot}{\cdot}).
}
\]

\item \textbf{$\beta$-Fourier coefficients of $v=\begin{bmatrix}1\\0\\ i\end{bmatrix}$.}
Since $\beta$ is orthonormal, coefficients are
\[
c_k=\ip{\beta_k}{v}\qquad (k=1,2,3).
\]
First compute
\[
Hv=
\begin{bmatrix}
2 & 1+i & 0\\
1-i & 3 & 2\\
0 & 2 & 4
\end{bmatrix}
\begin{bmatrix}1\\0\\ i\end{bmatrix}
=
\begin{bmatrix}
2\\
1+i\\
4i
\end{bmatrix}.
\]
Now:
\[
c_1=\ip{\beta_1}{v}
=\left(\frac{1}{\sqrt2}e_1\right)^\dagger Hv
=\frac{1}{\sqrt2}\,(2)=\sqrt2.
\]
Let $\alpha\df \frac{1+i}{2}$, so $u_2=e_2-\alpha e_1$ and $\overline{\alpha}=\frac{1-i}{2}$.
Then
\[
c_2=\ip{\beta_2}{v}
=\frac{1}{\sqrt2}\ip{u_2}{v}
=\frac{1}{\sqrt2}\,u_2^\dagger(Hv)
=\frac{1}{\sqrt2}\left((-\overline{\alpha})\cdot 2 + 1\cdot (1+i)\right)
=\frac{1}{\sqrt2}(2i)=\sqrt2\,i.
\]
Similarly, $u_3=\alpha e_1-e_2+e_3$, so
\[
c_3=\ip{\beta_3}{v}
=\frac{1}{\sqrt2}\,u_3^\dagger(Hv)
=\frac{1}{\sqrt2}\left(\overline{\alpha}\cdot 2 - (1+i) + 4i\right)
=\frac{1}{\sqrt2}(2i)=\sqrt2\,i.
\]
Therefore the Fourier coefficients are
\[
\boxed{\, (c_1,c_2,c_3)=\left(\sqrt2,\ \sqrt2\,i,\ \sqrt2\,i\right).\,}
\]
\end{enumerate}

\newpage
\subsubsection{Method 2: Matrix/Gram-matrix formulation}
The theorems and definitions used in this method are the following.
\begin{itemize}[leftmargin=*]
  \item \textbf{Hermitian positive definite matrix theorem.} If $H\in\C^{n,n}$ is Hermitian positive definite, then
  \[
  \ip{u}{v}\df u^\dagger Hv
  \]
  defines an inner product on $\C^{n,1}$ under the convention that the first slot is conjugate-linear and the second slot is linear.
  \item \textbf{Gram matrix definition.} Given a basis $(e_1,\dots,e_n)$ of an inner product space, its Gram matrix is
  $G=(\ip{e_i}{e_j})_{i,j}$. For the standard basis and the inner product $\ip{u}{v}=u^\dagger Hv$, this Gram matrix is exactly $H$.
\end{itemize}

\begin{enumerate}[label=(\alph*)]
\item Since
\[
H^\dagger=H,
\]
the matrix $H$ is Hermitian. Its leading principal minors are
\[
\Delta_1=2>0,
\]
\[
\Delta_2=\det\begin{bmatrix}2&1+i\\1-i&3\end{bmatrix}
=6-(1+i)(1-i)=6-2=4>0,
\]
and
\[
\Delta_3=\det(H)=8>0.
\]
By Sylvester's criterion for Hermitian matrices, $H$ is positive definite. Hence, by the Hermitian positive definite matrix theorem,
$u^\dagger Hv$ defines an inner product.
\[
\boxed{\,\ip{\cdot}{\cdot}\text{ is an inner product on }\C^{3,1}.\,}
\]

\item Let $e_1,e_2,e_3$ be the standard basis. Because the Gram matrix is $H$, we can read off
\[
\ip{e_i}{e_j}=H_{ij}.
\]
Thus
\[
\ip{e_1}{e_1}=2,
\qquad
\ip{e_1}{e_2}=1+i,
\qquad
\ip{e_1}{e_3}=0.
\]
Using the projection formula
\[
\proj_{u}(v)=\frac{\ip{u}{v}}{\ip{u}{u}}u,
\]
we get
\[
u_1=e_1,
\qquad
u_2=e_2-\frac{\ip{u_1}{e_2}}{\ip{u_1}{u_1}}u_1
=e_2-\frac{1+i}{2}e_1.
\]
Next, since
\[
\ip{u_1}{e_3}=0,
\qquad
\ip{u_2}{u_2}=2,
\qquad
\ip{u_2}{e_3}=2,
\]
we obtain
\[
u_3=e_3-\frac{\ip{u_1}{e_3}}{\ip{u_1}{u_1}}u_1
       -\frac{\ip{u_2}{e_3}}{\ip{u_2}{u_2}}u_2
=e_3-u_2
=\frac{1+i}{2}e_1-e_2+e_3.
\]
Therefore
\[
\boxed{\,
(u_1,u_2,u_3)=\left(e_1,\ e_2-\frac{1+i}{2}e_1,\ \frac{1+i}{2}e_1-e_2+e_3\right).
\,}
\]

\item The same Gram-matrix calculations give
\[
\ip{u_1}{u_1}=\ip{u_2}{u_2}=\ip{u_3}{u_3}=2.
\]
Hence
\[
\boxed{\,
\beta=\left(\frac{u_1}{\sqrt2},\frac{u_2}{\sqrt2},\frac{u_3}{\sqrt2}\right)
\,}
\]
is orthonormal, namely
\[
\boxed{
\beta=
\left(
\frac{1}{\sqrt2}\!\begin{bmatrix}1\\0\\0\end{bmatrix},\ 
\frac{1}{\sqrt2}\!\begin{bmatrix}-\frac{1+i}{2}\\1\\0\end{bmatrix},\ 
\frac{1}{\sqrt2}\!\begin{bmatrix}\frac{1+i}{2}\\-1\\1\end{bmatrix}
\right).
}
\]

\item For $v=(1,0,i)^{\mathsf T}$, first compute
\[
Hv=\begin{bmatrix}2\\1+i\\4i\end{bmatrix}.
\]
Since $\beta$ is orthonormal, the Fourier coefficients are
\[
c_k=\ip{\beta_k}{v}=\beta_k^\dagger Hv.
\]
Substituting the three vectors gives
\[
c_1=\sqrt2,
\qquad
c_2=\sqrt2\,i,
\qquad
c_3=\sqrt2\,i.
\]
Therefore
\[
\boxed{\,(c_1,c_2,c_3)=\left(\sqrt2,\sqrt2\,i,\sqrt2\,i\right).\,}
\]
\end{enumerate}

This method is efficient because it avoids recomputing each inner product from scratch: the entries of $H$ already encode all pairwise inner products of the standard basis vectors.

%====================================================
% QUESTION 5
%====================================================
\newpage
\section{Problem 5 (Double orthogonal complements)}
\subsection{Problem}
Let (i) $(V,\ip{\cdot}{\cdot})$ be an inner product space, and (ii) $W$ be a vector subspace of $V$.
\begin{enumerate}[label=(\alph*)]
\item Show that $W\subseteq W^{\perp\perp}$.
\item Show that if $V$ is finite-dimensional, then $W=W^{\perp\perp}$.
\item Show that if $V$ is infinite-dimensional, then it is not necessarily true that $W=W^{\perp\perp}$.
\end{enumerate}

\subsection{Solution}

\subsubsection{Method 1: Direct definition + dimension counting (finite-dimensional case)}
\begin{enumerate}[label=(\alph*)]
\item Let $w\in W$. To show $w\in W^{\perp\perp}$ we must show $\ip{w}{u}=0$ for all $u\in W^\perp$.
But $u\in W^\perp$ means precisely that $\ip{u}{w'}=0$ for all $w'\in W$; in particular $\ip{u}{w}=0$.
Using conjugate symmetry (or symmetry in the real case), $\ip{w}{u}=0$. Hence $w\in W^{\perp\perp}$.
So $W\subseteq W^{\perp\perp}$.

\[
\boxed{\,W\subseteq W^{\perp\perp}\,}
\]

\item Assume $\dim(V)<\infty$. A standard theorem gives
\[
\dim(W)+\dim(W^\perp)=\dim(V).
\]
Apply this to $W^\perp$:
\[
\dim(W^\perp)+\dim(W^{\perp\perp})=\dim(V).
\]
Subtracting yields $\dim(W^{\perp\perp})=\dim(W)$. Together with (a) and finite-dimensionality,
\[
W\subseteq W^{\perp\perp}\ \text{ and }\ \dim(W)=\dim(W^{\perp\perp})
\quad\Longrightarrow\quad
W=W^{\perp\perp}.
\]
\[
\boxed{\,\dim(V)<\infty\ \Rightarrow\ W=W^{\perp\perp}\,}
\]

\item In infinite dimensions, equality can fail. Let $V=\ell^2$ be the space of square-summable real sequences
(with the standard inner product $\ip{x}{y}=\sum_{k=1}^\infty x_k y_k$), and let
\[
W=c_{00}\df \{x\in\ell^2:\ x\text{ has finite support}\}.
\]
Then $W$ is a subspace of $V$ and is infinite-dimensional (but not all of $\ell^2$).

We claim $W^\perp=\{0\}$. Indeed, if $z\in W^\perp$, then for each standard basis vector $e_k\in c_{00}$ we have
\[
0=\ip{z}{e_k}=z_k,
\]
so all coordinates of $z$ are $0$, hence $z=0$.
Therefore $W^\perp=\{0\}$ and hence $W^{\perp\perp}=V$.

But $W\neq V$ (e.g.\ $(1,1/2,1/3,\dots)\in \ell^2$ but not finitely supported).
Thus $W\neq W^{\perp\perp}$.

\[
\boxed{\,\text{In infinite dimensions, }W\neq W^{\perp\perp}\text{ can occur (e.g.\ }W=c_{00}\subset \ell^2).\,}
\]
\end{enumerate}

\subsubsection{Method 2: Orthogonal decomposition (finite-dimensional case) + ``closure'' intuition}
\begin{enumerate}[label=(\alph*)]
\item Same as Method 1.

\item If $\dim(V)<\infty$, choose an orthonormal basis of $W$ and extend it to an orthonormal basis of $V$ by Gram--Schmidt:
\[
(\text{ONB of }W)\ \cup\ (\text{ONB of }W^\perp)=\text{ONB of }V.
\]
Then every $v\in V$ has a unique decomposition $v=w+u$ with $w\in W$, $u\in W^\perp$.
Now if $v\in W^{\perp\perp}$, then $\ip{v}{u}=0$ for all $u\in W^\perp$.
In particular, taking $u$ from the decomposition $v=w+u$ gives
\[
0=\ip{v}{u}=\ip{w+u}{u}=\ip{w}{u}+\ip{u}{u} = 0+\|u\|^2,
\]
so $u=0$ and hence $v=w\in W$. Thus $W^{\perp\perp}\subseteq W$, and with (a) we get equality.

\item In general (especially in infinite dimensions), $W^{\perp\perp}$ corresponds to the ``closed span'' of $W$ in the norm topology,
so if $W$ is not closed, strict inclusion may occur. The $\ell^2$ example above realizes this phenomenon concretely.
\end{enumerate}

%====================================================
% QUESTION 6
%====================================================
\newpage
\section{Problem 6 (Orthogonal complement and orthogonal decomposition in $\R^4$)}
\subsection{Problem}
Let
\[
V\df \{(w,x,y,z)\in\R^4\mid w-y+z=0\},
\]
which is a vector subspace of $\R^4$ (standard inner product $\cdot$).
\begin{enumerate}[label=(\alph*)]
\item Derive a basis for $V^\perp$.
\item Derive an orthonormal basis for $(V,\cdot)$.
\item Compute the unique $a\in V$ and $b\in V^\perp$ such that $(1,2,3,4)=a+b$.
\end{enumerate}

\subsection{Solution}

\subsubsection{Method 1: Solve linear constraints for $V^\perp$; then Gram--Schmidt inside $V$}
\begin{enumerate}[label=(\alph*)]
\item Parameterize $V$ by free variables $x,y,z$:
\[
w=y-z,\qquad (w,x,y,z)=(y-z,\ x,\ y,\ z)=x(0,1,0,0)+y(1,0,1,0)+z(-1,0,0,1).
\]
Let $u=(a,b,c,d)\in\R^4$. Then $u\in V^\perp$ iff $u\cdot v=0$ for all $v\in V$.
It suffices to impose orthogonality to the spanning vectors:
\[
u\cdot(0,1,0,0)=b=0,\qquad
u\cdot(1,0,1,0)=a+c=0,\qquad
u\cdot(-1,0,0,1)=-a+d=0.
\]
Thus $b=0$, $c=-a$, $d=a$, so
\[
V^\perp=\Span\{(1,0,-1,1)\}.
\]
\[
\boxed{\,V^\perp=\Span\{(1,0,-1,1)\}\,}.
\]

\item A convenient basis of $V$ from the parameterization is
\[
v_1=(0,1,0,0),\quad v_2=(1,0,1,0),\quad v_3=(-1,0,0,1).
\]
Apply Gram--Schmidt:

\underline{Step 1:} $e_1=v_1$ (already unit), so $e_1=(0,1,0,0)$.

\underline{Step 2:} $v_2\cdot e_1=0$, so $u_2=v_2$ and
\[
e_2=\frac{v_2}{\|v_2\|}=\frac{1}{\sqrt2}(1,0,1,0).
\]

\underline{Step 3:} Orthogonalize $v_3$ against $u_2=v_2$:
\[
v_3\cdot v_2=(-1)\cdot 1+0\cdot 0+0\cdot 1+1\cdot 0=-1,\qquad v_2\cdot v_2=2.
\]
So
\[
u_3=v_3-\proj_{v_2}(v_3)=v_3-\frac{-1}{2}v_2=v_3+\frac12 v_2=\left(-\frac12,0,\frac12,1\right).
\]
Equivalently, $2u_3=(-1,0,1,2)$, and $\|(-1,0,1,2)\|^2=6$. Hence
\[
e_3=\frac{1}{\sqrt6}(-1,0,1,2).
\]
Thus an orthonormal basis for $V$ is
\[
\boxed{\,
\left(
(0,1,0,0),\ \frac{1}{\sqrt2}(1,0,1,0),\ \frac{1}{\sqrt6}(-1,0,1,2)
\right).
\,}
\]

\item Let $p=(1,2,3,4)$. Since $V^\perp=\Span\{u\}$ with $u=(1,0,-1,1)$, we have
\[
b=\proj_{V^\perp}(p)=\frac{p\cdot u}{u\cdot u}u.
\]
Compute
\[
p\cdot u=1\cdot 1+2\cdot 0+3\cdot(-1)+4\cdot 1=2,\qquad u\cdot u=1+0+1+1=3.
\]
So
\[
b=\frac{2}{3}(1,0,-1,1)=\left(\frac23,0,-\frac23,\frac23\right).
\]
Then
\[
a=p-b=\left(1-\frac23,\ 2,\ 3+\frac23,\ 4-\frac23\right)=\left(\frac13,\ 2,\ \frac{11}{3},\ \frac{10}{3}\right).
\]
Check $a\in V$:
\[
w-y+z=\frac13-\frac{11}{3}+\frac{10}{3}=0.
\]
Therefore
\[
\boxed{\,
a=\left(\frac13,\ 2,\ \frac{11}{3},\ \frac{10}{3}\right)\in V,
\quad
b=\left(\frac23,0,-\frac23,\frac23\right)\in V^\perp,
\quad
(1,2,3,4)=a+b.
\,}
\]
\end{enumerate}

\subsubsection{Method 2: Matrix viewpoint (constraint normal vector)}
The subspace $V=\{x\in\R^4: n\cdot x=0\}$ where $n=(1,0,-1,1)$ (since $w-y+z=0$).
Hence $V^\perp=\Span\{n\}$ immediately. The orthogonal decomposition of $p$ is
\[
p=\left(p-\frac{p\cdot n}{n\cdot n}n\right)+\frac{p\cdot n}{n\cdot n}n,
\]
matching the $a,b$ above.

%====================================================
% QUESTION 7
%====================================================
\newpage
\section{Problem 7 (Diagonal matrices and Frobenius orthogonal complement)}
\subsection{Problem}
Let $n\in\mathbb{N}$, and consider the real inner product space $(\R^{n,n},\ip{\cdot}{\cdot}_F)$, where
\[
\ip{A}{B}_F\df \tr(A^{\mathsf T}B)=\sum_{i=1}^n\sum_{j=1}^n a_{ij}b_{ij}.
\]
Let
\[
D\df \{A\in\R^{n,n}\mid A\ \text{is diagonal}\}.
\]
Derive a basis for $D^\perp$.

\subsection{Solution}

\subsubsection{Method 1: Characterize $D^\perp$ by testing against diagonal matrices}
Let $X=(x_{ij})\in\R^{n,n}$. Then $X\in D^\perp$ iff $\ip{X}{\Delta}_F=0$ for all diagonal $\Delta=\diag(d_1,\dots,d_n)$.
But
\[
\ip{X}{\Delta}_F=\sum_{i=1}^n\sum_{j=1}^n x_{ij}\Delta_{ij}=\sum_{i=1}^n x_{ii}d_i.
\]
For this to be $0$ for all choices of $(d_1,\dots,d_n)\in\R^n$, we must have $x_{ii}=0$ for every $i$.
Thus
\[
D^\perp=\{X\in\R^{n,n}: x_{11}=x_{22}=\cdots=x_{nn}=0\},
\]
i.e.\ the subspace of matrices with zero diagonal.

Let $E_{ij}$ be the standard matrix units. Then $\{E_{ij}: i\neq j\}$ spans precisely the zero-diagonal matrices and is linearly independent.
Therefore
\[
\boxed{\,D^\perp=\Span\{E_{ij}: i\neq j\}\quad\text{and}\quad \{E_{ij}: i\neq j\}\ \text{is a basis of }D^\perp.\,}
\]

\subsubsection{Method 2: Orthogonal direct sum decomposition}
Every matrix $A\in\R^{n,n}$ decomposes uniquely as
\[
A=\underbrace{\diag(a_{11},\dots,a_{nn})}_{\in D}\ +\ \underbrace{\bigl(A-\diag(a_{11},\dots,a_{nn})\bigr)}_{\text{zero diagonal}}.
\]
Moreover, diagonal matrices are Frobenius-orthogonal to zero-diagonal matrices, since
\[
\ip{\diag(d_1,\dots,d_n)}{X}_F=\sum_{i=1}^n d_i x_{ii}=0
\]
whenever $x_{ii}=0$. Hence $\R^{n,n}=D\oplus D^\perp$ and $D^\perp$ is exactly the zero-diagonal subspace,
with basis $\{E_{ij}:i\neq j\}$.

% ============================================================
\section{Practice Question: Legendre Orthogonal Basis and Fourier Coefficients}

\subsection{Problem}
Let $\cP_2(\R)$ be equipped with the inner product
\[
  \ip{f}{g}\df \int_{-1}^{1} f(x)g(x)\,dx.
\]
Consider the standard polynomial list
\[
  1,\quad x,\quad x^2.
\]

\begin{enumerate}[label=(\alph*)]
  \item Verify that $\ip{\cdot}{\cdot}$ is an inner product on $\cP_2(\R)$.

  \item Obtain an orthonormal basis
  \[
    \beta=(\beta_0,\beta_1,\beta_2)
  \]
  for $\cP_2(\R)$.

  \item Let
  \[
    f(x)=2+3x+4x^2.
  \]
  Find the $\beta$-Fourier coefficients of $f$, i.e. find
  \[
    c_0=\ip{\beta_0}{f},\qquad
    c_1=\ip{\beta_1}{f},\qquad
    c_2=\ip{\beta_2}{f}.
  \]

  \item Hence express $f(x)$ as a linear combination of the orthonormal basis vectors in $\beta$.
\end{enumerate}

\newpage
\subsection{Solution}

\subsubsection{Method 1: Legendre Shortcut}

\begin{enumerate}[label=(\alph*)]
  \item We verify the inner product axioms.

  For $a,b\in\R$ and $f_1,f_2,g\in\cP_2(\R)$,
  \[
    \ip{af_1+bf_2}{g}
    =
    \int_{-1}^{1}(af_1(x)+bf_2(x))g(x)\,dx
  \]
  \[
    =
    a\int_{-1}^{1}f_1(x)g(x)\,dx
    +
    b\int_{-1}^{1}f_2(x)g(x)\,dx
    =
    a\ip{f_1}{g}+b\ip{f_2}{g}.
  \]
  Similarly, the inner product is linear in the second variable.

  Symmetry follows from commutativity of multiplication:
  \[
    \ip{f}{g}
    =
    \int_{-1}^{1}f(x)g(x)\,dx
    =
    \int_{-1}^{1}g(x)f(x)\,dx
    =
    \ip{g}{f}.
  \]

  For positive definiteness,
  \[
    \ip{f}{f}
    =
    \int_{-1}^{1}f(x)^2\,dx
    \ge 0.
  \]
  If
  \[
    \int_{-1}^{1}f(x)^2\,dx=0,
  \]
  then since $f(x)^2\ge0$ and $f$ is continuous, we must have
  \[
    f(x)^2=0
    \qquad
    \text{for all }x\in[-1,1].
  \]
  Hence $f(x)=0$ for all $x\in[-1,1]$. Since $f$ is a polynomial, this implies
  \[
    f=0.
  \]
  Therefore
  \[
    \boxed{\ip{\cdot}{\cdot}\text{ is an inner product on }\cP_2(\R).}
  \]

  \item For the inner product
  \[
    \ip{f}{g}=\int_{-1}^{1}f(x)g(x)\,dx,
  \]
  the Legendre polynomials give an orthogonal basis. Up to degree $2$, they are
  \[
    P_0(x)=1,
    \qquad
    P_1(x)=x,
    \qquad
    P_2(x)=\frac12(3x^2-1).
  \]
  They satisfy
  \[
    \ip{P_m}{P_n}=0
    \qquad(m\ne n),
  \]
  and
  \[
    \|P_n\|^2=\ip{P_n}{P_n}=\frac{2}{2n+1}.
  \]
  Hence
  \[
    \|P_0\|^2=2,
    \qquad
    \|P_1\|^2=\frac23,
    \qquad
    \|P_2\|^2=\frac25.
  \]
  Therefore the corresponding orthonormal basis is
  \[
    \beta_0=\frac{P_0}{\|P_0\|}
    =
    \frac1{\sqrt2},
  \]
  \[
    \beta_1=\frac{P_1}{\|P_1\|}
    =
    \sqrt{\frac32}x,
  \]
  and
  \[
    \beta_2=\frac{P_2}{\|P_2\|}
    =
    \sqrt{\frac52}\cdot \frac12(3x^2-1).
  \]
  Thus
  \[
    \boxed{
    \beta=
    \left(
    \frac1{\sqrt2},
    \sqrt{\frac32}x,
    \sqrt{\frac52}\cdot\frac12(3x^2-1)
    \right).}
  \]

  \item Since $\beta$ is orthonormal, the Fourier coefficients are
  \[
    c_i=\ip{\beta_i}{f}.
  \]

  First,
  \[
    c_0
    =
    \ip{\frac1{\sqrt2}}{2+3x+4x^2}
    =
    \frac1{\sqrt2}\int_{-1}^{1}(2+3x+4x^2)\,dx.
  \]
  Using
  \[
    \int_{-1}^{1}1\,dx=2,
    \qquad
    \int_{-1}^{1}x\,dx=0,
    \qquad
    \int_{-1}^{1}x^2\,dx=\frac23,
  \]
  we get
  \[
    c_0
    =
    \frac1{\sqrt2}\left(2\cdot2+3\cdot0+4\cdot\frac23\right)
    =
    \frac1{\sqrt2}\left(4+\frac83\right)
    =
    \frac{20}{3\sqrt2}.
  \]
  Equivalently,
  \[
    \boxed{c_0=\frac{10\sqrt2}{3}.}
  \]

  Next,
  \[
    c_1
    =
    \ip{\sqrt{\frac32}x}{2+3x+4x^2}
    =
    \sqrt{\frac32}\int_{-1}^{1}x(2+3x+4x^2)\,dx.
  \]
  Hence
  \[
    c_1
    =
    \sqrt{\frac32}\int_{-1}^{1}(2x+3x^2+4x^3)\,dx.
  \]
  The odd terms integrate to $0$, so
  \[
    c_1
    =
    \sqrt{\frac32}\left(3\cdot\frac23\right)
    =
    2\sqrt{\frac32}.
  \]
  Therefore
  \[
    \boxed{c_1=\sqrt6.}
  \]

  Finally,
  \[
    c_2
    =
    \ip{\sqrt{\frac52}\cdot\frac12(3x^2-1)}{2+3x+4x^2}.
  \]
  Thus
  \[
    c_2
    =
    \sqrt{\frac52}\int_{-1}^{1}
    \frac12(3x^2-1)(2+3x+4x^2)\,dx.
  \]
  Since $P_2$ is orthogonal to constants and to $x$, only the $4x^2$ term contributes:
  \[
    c_2
    =
    \sqrt{\frac52}\int_{-1}^{1}
    \frac12(3x^2-1)(4x^2)\,dx.
  \]
  So
  \[
    c_2
    =
    \sqrt{\frac52}\cdot 2\int_{-1}^{1}(3x^4-x^2)\,dx.
  \]
  Using
  \[
    \int_{-1}^{1}x^4\,dx=\frac25,
    \qquad
    \int_{-1}^{1}x^2\,dx=\frac23,
  \]
  we get
  \[
    c_2
    =
    \sqrt{\frac52}\cdot 2\left(3\cdot\frac25-\frac23\right)
    =
    \sqrt{\frac52}\cdot 2\left(\frac65-\frac23\right).
  \]
  Since
  \[
    \frac65-\frac23=\frac{18-10}{15}=\frac8{15},
  \]
  we obtain
  \[
    c_2
    =
    \sqrt{\frac52}\cdot\frac{16}{15}.
  \]
  Therefore
  \[
    \boxed{c_2=\frac{16}{15}\sqrt{\frac52}.}
  \]

  Hence the $\beta$-Fourier coefficients are
  \[
    \boxed{
    (c_0,c_1,c_2)
    =
    \left(
    \frac{10\sqrt2}{3},
    \sqrt6,
    \frac{16}{15}\sqrt{\frac52}
    \right).}
  \]

  \item Since $\beta$ is orthonormal,
  \[
    f=c_0\beta_0+c_1\beta_1+c_2\beta_2.
  \]
  Therefore
  \[
    \boxed{
    f(x)
    =
    \frac{10\sqrt2}{3}\left(\frac1{\sqrt2}\right)
    +
    \sqrt6\left(\sqrt{\frac32}x\right)
    +
    \frac{16}{15}\sqrt{\frac52}
    \left(\sqrt{\frac52}\cdot\frac12(3x^2-1)\right).}
  \]
  Simplifying,
  \[
    \frac{10\sqrt2}{3}\cdot\frac1{\sqrt2}
    =
    \frac{10}{3},
  \]
  \[
    \sqrt6\cdot\sqrt{\frac32}x
    =
    3x,
  \]
  and
  \[
    \frac{16}{15}\sqrt{\frac52}\cdot\sqrt{\frac52}\cdot\frac12(3x^2-1)
    =
    \frac{16}{15}\cdot\frac52\cdot\frac12(3x^2-1)
    =
    \frac43(3x^2-1).
  \]
  Hence
  \[
    f(x)
    =
    \frac{10}{3}+3x+\frac43(3x^2-1).
  \]
  Thus
  \[
    f(x)
    =
    \frac{10}{3}+3x+4x^2-\frac43
    =
    2+3x+4x^2,
  \]
  as required.
\end{enumerate}

\newpage
\subsubsection{Method 2: Faster Orthogonal Legendre Expansion}

Instead of using the orthonormal basis directly, we may first expand $f$ in the
orthogonal Legendre basis
\[
  P_0=1,
  \qquad
  P_1=x,
  \qquad
  P_2=\frac12(3x^2-1).
\]
Since
\[
  P_2=\frac12(3x^2-1),
\]
we have
\[
  3x^2=2P_2+1,
\]
so
\[
  x^2=\frac13P_0+\frac23P_2.
\]
Therefore
\[
  f(x)=2+3x+4x^2
  =
  2P_0+3P_1+4\left(\frac13P_0+\frac23P_2\right).
\]
Hence
\[
  f(x)
  =
  \frac{10}{3}P_0+3P_1+\frac83P_2.
\]
Thus, in the orthogonal Legendre basis,
\[
  \boxed{
  f(x)=\frac{10}{3}P_0(x)+3P_1(x)+\frac83P_2(x).}
\]

Now convert to the orthonormal basis
\[
  \beta_0=\frac{P_0}{\|P_0\|},
  \qquad
  \beta_1=\frac{P_1}{\|P_1\|},
  \qquad
  \beta_2=\frac{P_2}{\|P_2\|}.
\]
Since
\[
  \|P_0\|=\sqrt2,
  \qquad
  \|P_1\|=\sqrt{\frac23},
  \qquad
  \|P_2\|=\sqrt{\frac25},
\]
we have
\[
  P_0=\sqrt2\,\beta_0,
  \qquad
  P_1=\sqrt{\frac23}\,\beta_1,
  \qquad
  P_2=\sqrt{\frac25}\,\beta_2.
\]
Therefore
\[
  f
  =
  \frac{10}{3}\sqrt2\,\beta_0
  +
  3\sqrt{\frac23}\,\beta_1
  +
  \frac83\sqrt{\frac25}\,\beta_2.
\]
So the Fourier coefficients are
\[
  \boxed{
  c_0=\frac{10\sqrt2}{3},
  \qquad
  c_1=\sqrt6,
  \qquad
  c_2=\frac83\sqrt{\frac25}.}
\]
Since
\[
  \frac83\sqrt{\frac25}
  =
  \frac{16}{15}\sqrt{\frac52},
\]
this agrees with Method 1.

\end{document}